% =====================================================================
%  MACROS PROPIAS
%  Añade aquí tus comandos (\newcommand) para no repetir notación.
% =====================================================================

% --- Resumen / Abstract ----------------------------------------------
% Entorno para el "RESUMEN DEL PROYECTO" y el "ABSTRACT" que exige la
% normativa. Las figuras que se incluyan dentro se numeran aparte (1, 2…)
% y no aparecen en el índice de figuras.
%   #1: título (RESUMEN DEL PROYECTO / ABSTRACT)
\NewDocumentEnvironment{resumenproyecto}{m}{%
  \cleardoublepage
  \thispagestyle{plain}%
  \setcounter{figure}{0}%
  \renewcommand{\thefigure}{\arabic{figure}}%
  \renewcommand{\theHfigure}{#1.\arabic{figure}}%
  \captionsetup{list=false}%
  \begin{center}\large\bfseries #1\end{center}
  \vspace{1em}
}{%
  \setcounter{figure}{0}%
  \clearpage
}

% Apartado numerado del resumen (1. Introducción, 2. Metodología…)
\newcommand{\apartadoresumen}[1]{\par\medskip\noindent\textbf{#1}\par\smallskip\noindent}

% --- Notas de trabajo ---------------------------------------------------
% \pendiente{...} marca en rojo lo que falta por hacer. La CI cuenta
% cuántos quedan en cada compilación, así que conviene usarlo.
\newcommand{\pendiente}[1]{\textcolor{red}{\textbf{[PENDIENTE:} #1\textbf{]}}}
% \nota{...} para comentarios del director o del alumno en los borradores
\newcommand{\nota}[1]{\textcolor{blue}{\textit{[Nota: #1]}}}

% --- Notación (ejemplos) ---------------------------------------------------
\newcommand{\vect}[1]{\boldsymbol{#1}}
\newcommand{\dif}{\mathop{}\!\mathrm{d}}
